\documentclass{article}

% --- PACKAGES ---
\usepackage{fontspec}
\usepackage{geometry}
\usepackage[colorlinks=true, urlcolor=blue, linkcolor=blue]{hyperref}
\usepackage{graphicx}
\usepackage{longtable}
\usepackage{booktabs}

% --- GEOMETRY ---
\geometry{a4paper, margin=2.5cm}

% --- FONTS ---
% Set the main font for the document (body text, headings)
\setmainfont{Inter}[Path=./docs/fonts/inter-fonts/extras/otf/]
\setsansfont{Inter}[Path=./docs/fonts/inter-fonts/extras/otf/]

% Set the monospace font (for code blocks)
\setmonofont{JetBrains Mono}[Path=./docs/fonts/fonts/ttf/]

% --- PANDOC VARIABLES ---

% --- DOCUMENT ---
\begin{document}


  {
    \hypersetup{linkcolor=black}
    \tableofcontents
  }

\section{Planorama Business Plan}\label{planorama-business-plan}

\subsection{Inhaltsverzeichnis}\label{inhaltsverzeichnis}

\begin{enumerate}
\def\labelenumi{\arabic{enumi}.}
\tightlist
\item
  \hyperref[innovative-idee-planorama]{Innovative Idee: Planorama}
\item
  \hyperref[business-model-canvas-planorama]{Business Model Canvas:
  Planorama}
\item
  \hyperref[st-galler-business-model-navigator-planorama]{St.~Galler
  Business Model Navigator: Planorama}
\item
  \hyperref[design-sprint-planorama]{Design Sprint: Planorama}
\item
  \hyperref[innovatives-framework-nextjs]{Innovatives Framework:
  Next.js}
\item
  \hyperref[innovationen-bewerten-pestel-analyse-fuxfcr-planorama]{Innovationen
  bewerten: PESTEL-Analyse für Planorama}
\item
  \hyperref[stakeholder-analyse-planorama]{Stakeholder-Analyse:
  Planorama}
\item
  \hyperref[marktforschung-planorama]{Marktforschung: Planorama}
\item
  \hyperref[erfolgsmessung-planorama]{Erfolgsmessung: Planorama}
\item
  \hyperref[risikoanalyse-planorama]{Risikoanalyse: Planorama}
\item
  \hyperref[pruxe4sentationstechnik-pitch]{Präsentationstechnik: Pitch}
\end{enumerate}

\newpage

\section{Innovative Idee: Planorama}\label{innovative-idee-planorama}

\textbf{Planorama} ist eine innovative Plattform zur Urlaubsplanung, die
es Nutzern ermöglicht, Reisen kollaborativ zu organisieren, Inspiration
zu finden und alle logistischen Aspekte an einem Ort zu verwalten.

\subsection{Zweck}\label{zweck}

Planorama hilft Einzelpersonen und Gruppen, Urlaube effizient zu planen,
indem es Tools für die gemeinsame Planung, Aufgabenverteilung und
Inspiration bietet.

\subsection{Zielgruppe}\label{zielgruppe}

\begin{itemize}
\tightlist
\item
  Familien
\item
  Freundesgruppen
\item
  Paare
\item
  Einzelreisende
\end{itemize}

\subsection{Innovationsaspekt}\label{innovationsaspekt}

\begin{itemize}
\tightlist
\item
  Kollaborative Planung in Echtzeit
\item
  Integration von Inspiration, Buchung und Aufgabenmanagement
\item
  Intuitive Benutzeroberfläche
\item
  Personalisierte Empfehlungen basierend auf Nutzerpräferenzen
\end{itemize}

\newpage

\section{Business Model Canvas:
Planorama}\label{business-model-canvas-planorama}

\subsection{Kundensegmente}\label{kundensegmente}

\begin{itemize}
\tightlist
\item
  Familien, Freundesgruppen, Paare, Einzelreisende
\end{itemize}

\subsection{Wertangebote}\label{wertangebote}

\begin{itemize}
\tightlist
\item
  Kollaborative Urlaubsplanung
\item
  Inspiration und Empfehlungen
\item
  Zentrale Verwaltung aller Reisedaten
\end{itemize}

\subsection{Kanäle}\label{kanuxe4le}

\begin{itemize}
\tightlist
\item
  Web-App
\item
  Mobile App
\end{itemize}

\subsection{Kundenbeziehungen}\label{kundenbeziehungen}

\begin{itemize}
\tightlist
\item
  Personalisierter Support
\item
  Community-Features
\end{itemize}

\subsection{Einnahmequellen}\label{einnahmequellen}

\begin{itemize}
\tightlist
\item
  Freemium-Modell (Basis kostenlos, Premium-Features kostenpflichtig)
\item
  Affiliate-Links für Buchungen
\end{itemize}

\subsection{Schlüsselressourcen}\label{schluxfcsselressourcen}

\begin{itemize}
\tightlist
\item
  Plattform/Software
\item
  Partnerschaften mit Reiseanbietern
\end{itemize}

\subsection{Schlüsselaktivitäten}\label{schluxfcsselaktivituxe4ten}

\begin{itemize}
\tightlist
\item
  Entwicklung und Wartung der Plattform
\item
  Marketing
\item
  Kundenservice
\end{itemize}

\subsection{Schlüsselpartner}\label{schluxfcsselpartner}

\begin{itemize}
\tightlist
\item
  Reiseanbieter
\item
  Buchungsplattformen
\end{itemize}

\subsection{Kostenstruktur}\label{kostenstruktur}

\begin{itemize}
\tightlist
\item
  Entwicklungskosten
\item
  Marketing
\item
  Support
\end{itemize}

\begin{figure}
\centering
\pandocbounded{\includegraphics[keepaspectratio,alt={Business Model Canvas}]{docs/business-model-canvas-planorama.png}}
\caption{Business Model Canvas}
\end{figure}

\newpage

\section{St.~Galler Business Model Navigator:
Planorama}\label{st.-galler-business-model-navigator-planorama}

\subsection{Gewähltes Modell:
Freemium}\label{gewuxe4hltes-modell-freemium}

Planorama nutzt das Freemium-Modell, bei dem die Grundfunktionen
kostenlos angeboten werden und Premium-Features gegen Bezahlung
freigeschaltet werden können.

\subsubsection{Begründung}\label{begruxfcndung}

\begin{itemize}
\tightlist
\item
  Niedrige Einstiegshürde für neue Nutzer
\item
  Möglichkeit, durch Premium-Features zusätzliche Einnahmen zu
  generieren
\item
  Nutzer können die Plattform risikofrei testen und bei Gefallen
  upgraden
\end{itemize}

\begin{figure}
\centering
\pandocbounded{\includegraphics[keepaspectratio,alt={St.~Galler Business Model}]{docs/stgaller-business-model-planorama.png}}
\caption{St.~Galler Business Model}
\end{figure}

\newpage

\section{Design Sprint: Planorama}\label{design-sprint-planorama}

\subsection{Ziel}\label{ziel}

Die wichtigsten Annahmen und Funktionen von Planorama schnell
validieren.

\subsection{Ablauf}\label{ablauf}

\begin{enumerate}
\def\labelenumi{\arabic{enumi}.}
\tightlist
\item
  \textbf{Verstehen:} Zielgruppe und deren Bedürfnisse analysieren
\item
  \textbf{Skizzieren:} Lösungsansätze für die Urlaubsplanung entwerfen
\item
  \textbf{Entscheiden:} Beste Ideen auswählen
\item
  \textbf{Prototypen:} Interaktiven Prototypen der Plattform erstellen
\item
  \textbf{Testen:} Prototyp mit echten Nutzern testen und Feedback
  einholen
\end{enumerate}

\subsection{Ergebnis}\label{ergebnis}

\begin{itemize}
\tightlist
\item
  Erkenntnisse über Nutzerbedürfnisse
\item
  Validierung der wichtigsten Funktionen
\item
  Iterationsgrundlage für die Weiterentwicklung
\end{itemize}

\begin{figure}
\centering
\pandocbounded{\includegraphics[keepaspectratio,alt={Planorama Design}]{docs/planorma-design.png}}
\caption{Planorama Design}
\end{figure}

\newpage

\section{Innovatives Framework:
Next.js}\label{innovatives-framework-next.js}

\subsection{Auswahl}\label{auswahl}

Für die Umsetzung von Planorama wurde das Framework \textbf{Next.js}
gewählt.

\subsection{Begründung}\label{begruxfcndung-1}

\begin{itemize}
\tightlist
\item
  Moderne, performante Webentwicklung
\item
  Server-Side Rendering für schnelle Ladezeiten
\item
  Gute Integration mit React
\item
  Umfangreiche Community und viele Ressourcen
\item
  Skalierbarkeit für zukünftige Erweiterungen
\end{itemize}

\newpage

\section{Innovationen bewerten: PESTEL-Analyse für
Planorama}\label{innovationen-bewerten-pestel-analyse-fuxfcr-planorama}

\subsection{Politisch}\label{politisch}

\begin{itemize}
\tightlist
\item
  Datenschutzbestimmungen für Nutzerdaten
\end{itemize}

\subsection{Wirtschaftlich}\label{wirtschaftlich}

\begin{itemize}
\tightlist
\item
  Wachsender Markt für Online-Reiseplanung
\item
  Potenzial für Affiliate-Einnahmen
\end{itemize}

\subsection{Sozial}\label{sozial}

\begin{itemize}
\tightlist
\item
  Trend zu individuellen und Gruppenreisen
\item
  Steigende Nutzung digitaler Tools
\end{itemize}

\subsection{Technologisch}\label{technologisch}

\begin{itemize}
\tightlist
\item
  Fortschritte bei Web- und Mobile-Technologien
\end{itemize}

\subsection{Umwelt}\label{umwelt}

\begin{itemize}
\tightlist
\item
  Möglichkeit, nachhaltige Reiseoptionen zu integrieren
\end{itemize}

\subsection{Rechtlich}\label{rechtlich}

\begin{itemize}
\tightlist
\item
  Einhaltung von DSGVO und anderen Datenschutzgesetzen
\end{itemize}

\newpage

\section{Stakeholder-Analyse:
Planorama}\label{stakeholder-analyse-planorama}

\subsection{Methode:
Stakeholder-Matrix}\label{methode-stakeholder-matrix}

\begin{longtable}[]{@{}
  >{\raggedright\arraybackslash}p{(\linewidth - 6\tabcolsep) * \real{0.2625}}
  >{\raggedright\arraybackslash}p{(\linewidth - 6\tabcolsep) * \real{0.2375}}
  >{\raggedright\arraybackslash}p{(\linewidth - 6\tabcolsep) * \real{0.1750}}
  >{\raggedright\arraybackslash}p{(\linewidth - 6\tabcolsep) * \real{0.3250}}@{}}
\toprule\noalign{}
\begin{minipage}[b]{\linewidth}\raggedright
Stakeholder
\end{minipage} & \begin{minipage}[b]{\linewidth}\raggedright
Interesse
\end{minipage} & \begin{minipage}[b]{\linewidth}\raggedright
Einfluss
\end{minipage} & \begin{minipage}[b]{\linewidth}\raggedright
Massnahmen
\end{minipage} \\
\midrule\noalign{}
\endhead
\bottomrule\noalign{}
\endlastfoot
Nutzer & Hoch & Hoch & Feedback einholen, Support \\
Reiseanbieter & Mittel & Mittel & Partnerschaften, API \\
Investoren & Hoch & Hoch & Transparente Kommunikation \\
Entwicklerteam & Hoch & Hoch & Agile Entwicklung \\
Gesetzgeber & Niedrig & Mittel & Einhaltung der Gesetze \\
\end{longtable}

\subsection{Beispiel
Stakeholder-Kommunikation}\label{beispiel-stakeholder-kommunikation}

\textbf{Beispiel-E-Mail an Reiseanbieter:}

Betreff: Partnerschaftsanfrage -- Planorama, die neue Plattform für
Urlaubsplanung

Sehr geehrte Damen und Herren,

wir entwickeln mit Planorama eine innovative Plattform, die Gruppen und
Einzelpersonen bei der Urlaubsplanung unterstützt. Wir sind überzeugt,
dass eine Partnerschaft mit Ihnen für beide Seiten Mehrwert schaffen
kann -- z.B. durch die Integration Ihrer Angebote in unsere
Buchungsstrecke.

Gerne stellen wir Ihnen unser Konzept in einem kurzen Gespräch vor und
freuen uns auf Ihre Rückmeldung.

Mit freundlichen Grüßen\\
Das Planorama-Team

\newpage

\section{Marktforschung: Planorama}\label{marktforschung-planorama}

\subsection{Primärforschung}\label{primuxe4rforschung}

\begin{itemize}
\tightlist
\item
  Online-Umfrage unter potenziellen Nutzern zu Reisegewohnheiten und
  Planungsproblemen
\item
  Interviews mit Reisegruppen
\end{itemize}

\subsubsection{Beispiel-Umfragefragen}\label{beispiel-umfragefragen}

\begin{enumerate}
\def\labelenumi{\arabic{enumi}.}
\tightlist
\item
  Wie oft planen Sie Urlaubsreisen pro Jahr?
\item
  Mit wem planen Sie Ihre Reisen meistens? (Alleine, Familie, Freunde,
  Partner/in)
\item
  Welche Tools oder Apps nutzen Sie aktuell zur Urlaubsplanung?
\item
  Was sind die größten Herausforderungen bei der Planung von
  Gruppenreisen?
\item
  Wie wichtig sind Ihnen Inspirationen (z.B. Reisevorschläge,
  Aktivitäten) bei der Planung?
\item
  Welche Funktionen würden Sie sich von einer Urlaubsplanungs-App
  wünschen?
\item
  Wie bewerten Sie die Übersichtlichkeit und Benutzerfreundlichkeit
  Ihrer aktuellen Tools? (Skala 1-5)
\item
  Wären Sie bereit, für Premium-Funktionen einer
  Urlaubsplanungsplattform zu bezahlen?
\item
  Wie wichtig ist Ihnen der Datenschutz bei der Nutzung solcher
  Plattformen?
\item
  Haben Sie schon einmal eine Reise abgebrochen oder umgeplant, weil die
  Organisation zu kompliziert war?
\end{enumerate}

\subsection{Sekundärforschung}\label{sekunduxe4rforschung}

\begin{itemize}
\tightlist
\item
  Analyse von Marktstudien zum Online-Reisemarkt
\item
  Recherche zu bestehenden Lösungen und deren Schwächen
\end{itemize}

\newpage

\section{Erfolgsmessung: Planorama}\label{erfolgsmessung-planorama}

\subsection{Methode: Leistungskennzahlen
(KPIs)}\label{methode-leistungskennzahlen-kpis}

\begin{itemize}
\tightlist
\item
  Anzahl aktiver Nutzer
\item
  Anzahl erstellter Reisepläne
\item
  Nutzerzufriedenheit (z.B. durch Umfragen)
\item
  Conversion-Rate von kostenlosen zu zahlenden Nutzern
\end{itemize}

\subsubsection{Beispiel-KPIs}\label{beispiel-kpis}

\begin{itemize}
\tightlist
\item
  1.000 aktive Nutzer im ersten Quartal
\item
  100 erstellte Gruppenreisen pro Monat
\item
  80\% positive Bewertungen im App Store
\item
  10\% Conversion-Rate von Free zu Premium
\end{itemize}

\newpage

\section{Risikoanalyse: Planorama}\label{risikoanalyse-planorama}

\subsection{Methode: SWOT-Analyse}\label{methode-swot-analyse}

\subsubsection{Stärken}\label{stuxe4rken}

\begin{itemize}
\tightlist
\item
  Innovative, kollaborative Plattform
\item
  Einfache Bedienung
\item
  Kollaborative Planung in Echtzeit
\item
  Integration von Inspiration, Buchung und Aufgabenmanagement
\end{itemize}

\subsubsection{Schwächen}\label{schwuxe4chen}

\begin{itemize}
\tightlist
\item
  Abhängigkeit von Partnern
\item
  Markteintrittsbarrieren
\item
  Noch geringe Markenbekanntheit
\item
  Abhängigkeit von Partnerintegrationen
\end{itemize}

\subsubsection{Chancen}\label{chancen}

\begin{itemize}
\tightlist
\item
  Wachsende Nachfrage nach digitalen Planungstools
\item
  Erweiterung um neue Features
\item
  Kooperationen mit Reiseanbietern
\item
  Trend zu individuellen Gruppenreisen
\end{itemize}

\subsubsection{Risiken}\label{risiken}

\begin{itemize}
\tightlist
\item
  Starke Konkurrenz
\item
  Datenschutzanforderungen
\item
  Datenschutzprobleme könnten Vertrauen beeinträchtigen
\item
  Konkurrenz durch etablierte Plattformen
\end{itemize}

\newpage

\section{Präsentationstechnik:
Pitch}\label{pruxe4sentationstechnik-pitch}

\subsection{Auswahl}\label{auswahl-1}

Für die Präsentation von Planorama wird die Pitch-Technik gewählt.

\subsection{Beispiel-Pitch}\label{beispiel-pitch}

„Stellen Sie sich vor, Sie möchten mit Freunden oder Familie eine Reise
planen -- doch die Organisation wird schnell chaotisch: endlose
Chatverläufe, unübersichtliche Listen, vergessene Aufgaben.

Planorama löst dieses Problem! Unsere Plattform vereint Inspiration,
Planung und Aufgabenmanagement an einem Ort. Sie können gemeinsam
Reiseziele entdecken, Aufgaben verteilen, Unterkünfte buchen und alles
übersichtlich im Team organisieren -- egal ob am Desktop oder mobil.

Mit Planorama wird die Urlaubsplanung endlich einfach, transparent und
macht Spaß. Testen Sie kostenlos die Basisversion -- und entdecken Sie,
wie entspannt Reiseplanung sein kann!{\kern0pt}``

\subsection{Begründung}\label{begruxfcndung-2}

\begin{itemize}
\tightlist
\item
  Kompakte und überzeugende Darstellung der Idee
\item
  Fokus auf die wichtigsten Vorteile und Alleinstellungsmerkmale
\item
  Geeignet für Investoren, Partner und Nutzer
\end{itemize}

\end{document} 